\section{Supplemental Results}

\subsection{Experiment Overviews}

Below are overviews of the experiments performed. The colored regions are test case sequences. Times are PST. 

\begin{figure}[h]
    \includegraphics[width=0.8\linewidth]{\repodir/experiment/analysis/auto/output/figures/tc_plot_2023-04-07.png} 
    \caption{Experiment Overview for 2023-04-07 }
    \label{fig:SI_expt_overview_2023-04-07}
\end{figure}


\begin{figure}[h]
    \includegraphics[width=0.8\linewidth]{\repodir/experiment/analysis/auto/output/figures/tc_plot_2023-05-12.png} 
    \caption{Experiment Overview for 2023-05-12 }
    \label{fig:SI_expt_overview_2023-05-12}
\end{figure}


\begin{figure}[h]
    \includegraphics[width=0.8\linewidth]{\repodir/experiment/analysis/auto/output/figures/tc_plot_2023-05-18.png} 
    \caption{Experiment Overview for 2023-05-18}
    \label{fig:SI_expt_overview_2023-05-18}
\end{figure}

\begin{figure}[h]
    \includegraphics[width=0.8\linewidth]{\repodir/experiment/analysis/auto/output/figures/tc_plot_2023-05-24.png} 
    \caption{Experiment Overview for 2023-05-24}
    \label{fig:SI_expt_overview_2023-05-24}
\end{figure}

\clearpage

\subsection{Equivalence Ratio Dependence}

On 2023-05-24, we performed measurements with equivalence ratio varying from the standard value of 0.8. 

\begin{itemize}
    \item 516 - a position sweep at fixed equivalence ratio 0.6
    \item 5x3\_pos - a position and equivalence ratio sweep at 5 equivalence ratios and 0.1\% potassium mass fraction. The lower potassium mass fraction allowed for longer times on each test case and therefore AES nK measurements. 
    \item 5x6\_pos - a position and equivalence ratio sweep at 5 equivalence ratios and 1\% potassium mass fraction. The high potassium mass fraction meant short times on each test case with only enough time for MWS measurements. 
\end{itemize}

\subsubsection{5x3 and 5x6 Position Dependence Test Case Sequences}

This section contains additional measurements for the 5x3 and 5x6 test case sequences. These were 'two dimensional' parameter sweeps of both stage position and equivalence ratio. See Figure \ref{fig:SI_expt_overview_2023-05-24} for how the sequence was performed. 

Figure \ref{fig:SI_5x3_pos_nK_m3} shows the AES measurements for the 5x3\_pos test case. Note the calibration of the AES signal for this test case exhibited significant calibration drift (see \ref{fig:SI_5x3_pos_alpha}), and systematic error may be present in the nK values. 

\begin{figure}[]
\centering
\includegraphics[width=0.8\textwidth]{\repodir/experiment/analysis/tcs/output/figures/5x3_pos_nK_m3.png}
\caption{AES measurements 5x3\_pos}
\label{fig:SI_5x3_pos_nK_m3}
\end{figure}

\begin{figure}[]
\centering
\includegraphics[width=0.8\textwidth]{\repodir/experiment/analysis/tcs/output/figures/5x3_pos_alpha.png}
\caption{absorption data for AES measurements 5x3\_pos. The absorption offset at high wavelength is due to drift in calibration.}
\label{fig:SI_5x3_pos_alpha}
\end{figure}


\begin{figure}[]
\centering
\includegraphics[width=0.8\textwidth]{\repodir/experiment/analysis/tcs/output/figures/5x3_pos_mws_KOH.png}
\caption{MWS measurements 5x3\_pos. CFD centerline values for KOH are also shown.}
\label{fig:SI_5x3_pos_mws_KOH}
\end{figure}


\begin{figure}[]
\centering
\includegraphics[width=0.8\textwidth]{\repodir/experiment/analysis/tcs/output/figures/5x6_pos_mws_KOH.png}
\caption{MWS measurements 5x6\_pos. CFD centerline values for KOH are also shown.}
\label{fig:SI_5x6_pos_mws_KOH}
\end{figure}


\begin{figure}[]
\centering
\includegraphics[width=0.8\textwidth]{\repodir/experiment/analysis/tcs/output/figures/5x6_pos_mws_AStime_phi.png}
\caption{MWS AS time profiles for 5x6\_pos.}
\label{fig:SI_5x6_pos_mws_AStime_phi}
\end{figure}


\clearpage
\subsubsection{516 Position Dependence Test Case Sequence}

This section contains additional measurements for the 516 test case. This test case was a position sweep at fixed nominal equivalence ratio 0.6. The higher seed rate and longer time on each test case allows for more accurate determination of AS decay times than 5x3\_pos and 5x6\_pos test case sequences. We could also perform AES measurements which are presented in the main text. 



TODO: why is coordinate data 0.65? 

TODO: this should probably be in the main text.

\begin{figure}[]
\centering
\includegraphics[width=0.8\textwidth]{\repodir/experiment/analysis/tcs/output/figures/516_mws_decaytime.png}
\caption{Fitted decay time for MWS measurements 516. The fits leading to these time constants are shown in Figure \ref{fig:516_mws_AS_fit}.}
\label{fig:SI_516_mws_decaytime}
\end{figure}

\begin{figure}[]
\centering
\includegraphics[width=0.8\textwidth]{\repodir/experiment/analysis/tcs/output/figures/516_mws_AS_fit.png}
\caption{AS time profiles and exponential fits for MWS measurements 516.}
\label{fig:516_mws_AS_fit}
\end{figure}


\clearpage
\subsection{ICCD Camera Measurements}

This section contains additional measurements by the ICCD camera. Figure \ref{fig:SI_536_iccd} shows ICCD data taken at the 536 test case at the nominal position of 178 mm. Figure \ref{fig:SI_536_iccd}C shows that the photolumensence decay is very short compared the the mws decay. In previous experiments, we have investigated the time resolved spectral decay of the photolumensence and observed that it is entirely due to the decay of the potassium D1 line (TODO: check), and the 10ns decay profile closely matches the time profile of the laser pulse. 

TODO: add old measurements with spectrograph
TODO: calibrate xy

\begin{figure}[]
\centering
\includegraphics[width=0.8\textwidth]{\repodir/final/figures/SI/output/ICCD_SI1.png}
\caption{ICCD camera images taken at 1\% potassium mass fraction in the Run 53x on 2023-05-24. During this test case the camera acquired 15 gate delay sweeps, which are averaged for these images. A) Images are taken at different gate delays. The laser is position at the nominal position. B) A zoom in of the norminal position. The image is averaged across gate delays C) The average of the region of interest used in B across the spatial dimension}
\label{fig:SI_536_iccd}
\end{figure}


\subsection{Laser power dependence}

Figure \ref{fig:SI_power_dependence} shows the MWS signal as a function of the laser power. We observe that the AS signal increases super-linearly with the laser power. For high laser powers, the AS decay exhibits a steeper slope within the first few $\mu s$ after the laser pulse, indicating second-order effects. The exponential decay time constant of each power is similar, indicating a relatively power-independent first-order recombination process (SI).

\begin{figure}[h]
    \includegraphics[width=0.8\linewidth]{\repodir/final/figures/SI/output/Power_MWS.png} 
    \caption{Power dependence}
    \label{fig:SI_power_dependence}
\end{figure}

\begin{figure}[h]
    \includegraphics[width=0.8\linewidth]{\repodir/experiment/analysis/tcs/output/figures/fit_mws_dnedt_power.png} 
    \caption{Fit parameters for power dependence}
    \label{fig:SI_power_dependence_fit}
\end{figure}

\clearpage

\subsection{Other Supplemental Results}

Temporary section for supplemental figures


\begin{figure}[h]
    \includegraphics[width=0.8\linewidth]{\repodir/final/analysis/output/figures/krm_cfd_pos.png} 
    \caption{Power dependence}
    \label{fig:SI_krm_cfd_pos.png}
\end{figure}


\subsubsection{UV Absorption Cross Sections}

We compiled data on UV absorption cross sections to determine potential ionization pathways. Data was collected within the relevant wavelength range for potassium photoionization (190-285 nm). The consolidated data is show in \ref{fig:SI_UV_abscs}. TODO: replace citations. 

For potassium we use compiled photoionization cross sections (Verner) for solar and blackbody radiation.{[}17{]} Regarding the other constituent gasses, we first rule out absorption by the atomic species H and O through the NIST database.{[}18{]} The UV absorption of the hydrogen molecule requires light with wavelength shorter than about 111 nm.{[}19{]} The OH radical similarly does not absorb UV light until wavelengths shorter than about 170 nm.{[}20{]} empirical formulas for the wavelength and temperature dependent UV cross sections exist for CO2{[}21{]}, H2O,{[}22{]} and O2.{[}22{]} KOH is dissociated by UV light in the 200-300 nm range and consequently has a relatively large absorption cross section. We extract the KOH absorption cross section data from Weng et all and extrapolate it in the range 190-200 nm.{[}23{]} Weng et al. observe a weak difference in absorption cross section for different flame conditions at 1400 and 1800 K so we do not include a temperature dependence to the KOH cross section. CO does absorb very weakly Cameron bands in the range of 170 -- 210 nm. We include CO absorption through the room temperature absorption measurements performed by Thompson et al.{[}24{]}


\begin{figure}[h]
    \includegraphics[width=0.8\linewidth]{\repodir/modeling/viability/figure_panels/output/abscs.png} 
    \caption{Optical absorption cross sections of species with significant cross section and UV absorption in the relevant range. Left: Fixed temperature wavelength dependence. Right: Fixed wavelength temperature dependence}
    \label{fig:SI_UV_abscs}
\end{figure}

